\section{Introduction}
\label{sec:intro}

We set out to build a spam filter that is judged by what its mistakes cost, not by how many it
makes. For every email the filter takes one of three actions: deliver it to the inbox, send it to the
spam folder, or hold it for a person to inspect. The two errors are very different. A spam in the
inbox is a small nuisance; a real email in the spam folder may never be seen. We priced them as
follows:
\begin{center}
\small
\begin{tabular}{lcc}
\toprule
Action & real email ($y=0$) & spam ($y=1$) \\
\midrule
deliver       & $0$           & $\cfn=1$ \\
spam folder   & $\cfp=10$     & $0$ \\
inspect       & $\ci=0.3$     & $\ci=0.3$ \\
\bottomrule
\end{tabular}
\end{center}
Everything we report is the \emph{realized loss}: the average of these costs over held-out emails.

The report follows the project in the order we did it. We first built a shared pipeline for data,
decisions and scoring (Section~\ref{sec:setup}) and worked out which action is cheapest for a given
spam probability (Section~\ref{sec:rule}). That rule needs a probability, so we wrote two classifiers
(Section~\ref{sec:models}) and tested the rule against the usual accuracy rule as the two error costs
move apart (Section~\ref{sec:sweep}). The results showed that good decisions need trustworthy
probabilities, which we checked next (Section~\ref{sec:calib}). We then asked what happens when the
share of spam changes after training (Section~\ref{sec:shift}), and finally how much inspection is
worth and at what price it stops paying (Section~\ref{sec:inspect}). All models are written from scratch in Python and NumPy.
